\documentclass[11pt]{article}
\usepackage[margin=1in]{geometry}
\usepackage[T1]{fontenc}
\usepackage[utf8]{inputenc}
\usepackage{amsmath,amssymb,amsthm}
\usepackage{enumitem}

\title{ICPC World Finals 2024\\D. Doubles Horseback Wrestling}
\author{}
\date{}

\begin{document}
\maketitle

\section*{Problem Summary}

Riders $i$ and $j$ can form a valid pair if there exist ratings
$r_i \in [l_i, u_i]$ and $r_j \in [l_j, u_j]$ such that $r_i + r_j = s$.
Equivalently,
\[
    l_i + l_j \le s \le u_i + u_j.
\]
We must output a maximum-cardinality matching in this graph.

\section*{Key Observations}

\begin{itemize}[leftmargin=*]
    \item Sort riders by increasing $l$.
    \item For a fixed rider $i$, any partner $j$ with smaller or equal $l$ must satisfy
    $l_j \le s - l_i$ and $u_j \ge s - u_i$.
    \item If we process riders from large $l$ down to small $l$, then when rider $i$ is current, every still
    unprocessed rider with $l_j \le s - l_i$ is already known and can be placed in a pool.
    \item Among feasible partners in that pool, it is optimal to use the one with the smallest possible $u$.
    A larger $u$ can only help future riders more than a smaller $u$ can.
    \item If rider $i$ cannot be matched when it is processed, then it can never be matched later: all later
    riders have even smaller $l$, so rider $i$ would still need a partner from the same pool, and the lower
    bound on $u_j$ would only get harder to satisfy for rider $i$ itself.
\end{itemize}

\section*{Algorithm}

\begin{enumerate}[leftmargin=*]
    \item Sort riders by increasing $l$.
    \item Scan the sorted list from right to left.
    \item Maintain a multiset of still-unmatched riders to the left of the current position whose lower bound
    already satisfies $l_j \le s - l_i$. Store them ordered by $u_j$.
    \item For the current rider $i$, find in the multiset the rider with smallest $u_j \ge s - u_i$.
    \item If such a rider exists, match them and remove that rider from the multiset. Otherwise leave $i$
    unmatched forever.
\end{enumerate}

\section*{Correctness Proof}

We prove that the algorithm outputs a maximum matching.

\paragraph{Lemma 1.}
When rider $i$ is processed, the multiset contains exactly the unmatched riders $j$ to the left of $i$
such that $l_j \le s - l_i$.

\paragraph{Proof.}
We scan riders in decreasing order of $l_i$. A pointer grows over the left prefix and inserts every rider
whose $l_j$ has become small enough to satisfy $l_j \le s - l_i$. Since $s - l_i$ only increases during the
scan, every such rider is inserted exactly once. We remove riders only when they are matched or when a
former pool rider becomes the current rider itself. Therefore the multiset contains exactly the still-unmatched
eligible riders to the left. \qed

\paragraph{Lemma 2.}
If rider $i$ is matched, choosing among all feasible partners the one with minimum $u_j$ cannot decrease
the size of an optimal completion.

\paragraph{Proof.}
Let $j$ be the feasible partner with minimum $u_j$, and suppose some optimal solution matches $i$ with
another feasible rider $k$. Since both are feasible for $i$, we have
\[
u_j \le u_k, \qquad l_j \le s - l_i, \qquad u_j \ge s - u_i.
\]
Replace pair $(i,k)$ by $(i,j)$. Any other rider that could be matched with $j$ in the optimal solution also
has lower-bound requirement at most the requirement used for $i$ at this moment, and using the smaller
$u_j$ instead of $u_k$ only frees the larger $u_k$ for later riders. So we can swap $j$ into $i$'s pair and
use $k$ wherever $j$ was used, or leave $k$ unmatched if $j$ was unused. The number of pairs does not
decrease. \qed

\paragraph{Lemma 3.}
If rider $i$ has no feasible partner when processed, then there exists an optimal matching in which $i$ is
unmatched.

\paragraph{Proof.}
By Lemma 1, every possible partner of $i$ among still-unprocessed riders is already in the multiset. If no
feasible rider there has $u_j \ge s - u_i$, then no rider to the left can pair with $i$. Riders to the right have
already been processed and, in any pair containing $i$, rider $i$ would be the larger-$l$ endpoint. Hence
all possible pairings involving $i$ have already been considered, and none is feasible. Therefore $i$ can be
left unmatched in some optimal solution. \qed

\paragraph{Theorem.}
The algorithm outputs a maximum-cardinality set of disjoint valid pairs.

\paragraph{Proof.}
Consider the riders in the order processed by the algorithm. By Lemma 2, whenever the algorithm matches
the current rider, there exists an optimal solution consistent with that choice. By Lemma 3, whenever the
algorithm leaves the current rider unmatched, there exists an optimal solution consistent with that choice as
well. Inductively, after every step there remains an optimal solution consistent with all choices already made
by the algorithm. After the final step, the algorithm itself is such an optimal solution. \qed

\section*{Complexity Analysis}

Sorting costs $O(n \log n)$. Every rider is inserted into and removed from the multiset at most once, and
every query is one \texttt{lower\_bound}, so the total running time is $O(n \log n)$. The memory usage is
$O(n)$.

\section*{Implementation Notes}

\begin{itemize}[leftmargin=*]
    \item The multiset stores pairs $(u_j,\text{position})$ so duplicates in $u_j$ are handled naturally.
    \item A rider may enter the multiset before becoming the current rider later; in that case we remove its
    own entry before processing it.
    \item The output may list the pairs in any order.
\end{itemize}

\end{document}
